\section{Limitations and Broader Impact}
\label{sec:limitations}

\subsection{Limitations}

\paragraph{Scope.}
AmphionASR v3 is optimised for Mandarin and English, with a small
amount of Mandarin--English code-switching data. We do not claim broad
multilingual coverage and we do not benchmark on languages outside
those reported in Section~\ref{sec:experiments}.

\paragraph{Open multi-talker benchmarks.}
The v3 target-speaker path is substantially weaker on the public
LibriMix test sets ($65.83$\% WER on Libri2Mix, $87.06$\% on Libri3Mix)
than on our internal multi-talker set ($14.29$\% MER). The likely
causes -- training/test mismatch in reference-utterance distribution,
limited LibriMix-style data in Stage~3 -- are documented in
Section~\ref{sec:tsasr}, but until the next training iteration the
public multi-talker numbers should be read as a known weakness, not a
final result.

\paragraph{No explicit noise-robustness training (yet).}
Section~\ref{sec:noise} is deliberately short: this release does not
include dedicated noise-augmented training. The released model inherits
Whisper's encoder-level robustness, but is not benchmarked against
Qwen3-ASR-1.7B \cite{qwen2025qwen3asr} on noisy / far-field audio in
this report.

\paragraph{Inherited LLM failure modes.}
As discussed in Section~\ref{sec:analysis}, the model can repeat on
long silences, drift between languages on code-switched content, and
over-bias on prompt-supplied hotwords when a term is not actually
spoken. These are well-known characteristics of LLM-based ASR and we
are explicit that AmphionASR shares them.

\paragraph{Latency and streaming.}
The v3 release is offline-batch, served via vLLM. Streaming inference
and a measured first-token-latency report are out of scope for this
version.

\subsection{Broader Impact}

\paragraph{Surveillance risk.}
A target-speaker recognizer that operates from a short reference
utterance lowers the technical barrier to selectively transcribing a
designated person in multi-speaker audio. This capability has
legitimate uses (transcribing one's own voice in a meeting,
accessibility) but can also be misused for unconsented surveillance.
The released model card will state that the intended use is
user-initiated transcription with the consent of all recorded parties
and will explicitly call out unconsented surveillance as out-of-scope.

\paragraph{Bias.}
% TODO: insert per-accent / per-dialect WER breakdown when the
% KeSpeech-by-dialect split is added to the eval pack.
We have not yet measured per-accent WER systematically. KeSpeech is the
natural starting point for Mandarin dialectal coverage and will be
broken out in a follow-up.

\paragraph{Data provenance.}
% TODO: confirm provenance and licensing for every internal data source
% before public release of weights or checkpoints.
Public sources are used under their stated licenses
(see Table~\ref{tab:data}). Internal sources are subject to internal
data-governance review prior to any model-weight release.

\paragraph{LLM-assisted writing.}
Portions of this report were drafted with the assistance of large
language models. All technical claims, numbers, and citations were
verified by the authors against the corresponding primary sources.
% TODO: maintain ack/llm-usage.md as new sections are written and
% reflect the final disclosure here.
